% ==================== 七、AI 辅助设计的使用记录 ====================
\section{AI 辅助设计的使用记录}

\subsection{记录方式与口径}

本章记录本设计中五次实际使用大模型的过程。记录遵守两条自我约束：

\begin{itemize}
  \item \textbf{只记真实发生的事}——AI 给了什么、我判断了什么、最后改成了什么。没有发生过的「人工修改」一律不写，包括那些看起来更漂亮的例子。
  \item \textbf{缺点写具体到错在哪一句}。只写「AI 有幻觉、需人工核对」没有信息量；要写清它倾向了什么、为什么这个倾向在本作业里站不住。
\end{itemize}

五次使用覆盖了本设计的不同环节，其分布本身就是对「AI 能做什么、不能做什么」的答案：

\begin{table}[H]
\centering
\caption{五次使用记录概览}
\tabfont
\begin{tabularx}{\textwidth}{@{}L{1.3cm}L{3.0cm}L{3.4cm}X@{}}
\toprule
编号 & 使用场景 & 产出物 & 暴露的主要缺陷 \\
\midrule
U-01 & 模块划分维度的比较 & 模块划分方案与模块清单 & 对「交付约束」（每模块需有界面）不敏感 \\
U-02 & 下单链路的依赖推导 & 组件图的调用依赖 & 举出的具体例子本身是错的 \\
U-03 & 五处职责与接口分歧 & 职责表、组件图画法 & 在「好看」与「合规」冲突时选好看 \\
U-04 & 19 屏跨屏数据一致性核对 & 修正后的 19 屏界面 & 单屏自洽、跨屏互斥；补屏另起一套事实 \\
U-05 & 报告自身的图表与引用自检 & 修正后的章节与编号 & 引用关系一旦篇幅变长即失效 \\
\bottomrule
\end{tabularx}
\end{table}

\subsection{U-01：模块划分方式的比较与取舍}

\textbf{使用场景。} 作业要求「写明软件的模块划分」，但实验一的 36 条需求是平铺编号的，本身没有模块结构，模块必须重新构造。这一步没有唯一正确答案，需要先知道有哪几种切法。

\textbf{输入与生成结果。} 输入实验一的完整需求清单与实验二的要求，AI 给出三种切法：A 按角色切（顾客 / 前台 / 店长 / 平台）；B 按业务能力切（每模块负责一类业务，多角色共用）；C 按订单生命周期切（搜索 $\rightarrow$ 下单 $\rightarrow$ 支付 $\rightarrow$ 入住 $\rightarrow$ 退房 $\rightarrow$ 售后）。它推荐 B，并给出十个模块草案：账号与权限、门店与房源、定价、库存、预订与订单、支付与退款、房态与工单、评价、审计与数据保留、通知。

\textbf{优缺点分析。} 长处是\textbf{一次给出多个维度并说明各自代价}，比只给一个方案更容易做判断；把 36 条需求批量归类很省时，人工逐条归类容易漏。不足是\textbf{它对作业的交付约束不敏感}：作业要求「为每个功能模块设计 UI 界面」，而「通知」模块没有任何界面——短信与推送发生在业务事件背后，用户看不到它。AI 给的是一个技术上漂亮、在作业里却站不住的划分。此外，模块规模是否失衡需要人自己数一遍条数才看得出来，AI 不主动提醒：评价、通知、支付各只对应 1 条需求，预订与订单对应 8 条。

\textbf{修改优化后的结果。} 采用 B，但做了两处调整：① 通知模块按「无界面」的理由并入预订与订单，模块数从 10 降到 9；② REQ-009（取消与退款）原本整条挂在预订模块，导致「退款执行」归谁不清楚，把退款执行划到支付与退款，取消规则仍留在预订。另确认五条横切约束不挂单个模块：REQ-015 性能、REQ-017 可靠性、REQ-018 易用性、REQ-019 范围、REQ-027 输入校验。

\textbf{一处必须交代的出入。} 以上记录发生在模块划分的阶段，当时的结论是 9 个模块。报告第 2 章与组件图用的是 \textbf{10 个模块}——成稿时通知被重新独立为 M10，理由有三：其一，它在组件图上以「事件订阅者」的身份出现，与调用依赖的形态不同，合并进 M5 会让 M5 同时承担业务与通知两种职责；其二，REQ-029 要求通知失败进入独立失败队列并重试，这个失败语义不能污染订单事务；其三，通知的站内消息在界面上是可查的（消息记录与重试状态），因此「无界面」这个最初的排除理由在重新审视时不成立。需要说明的是，写这段话时第三条还只是一个判断——直到第 6 章补出 17 屏消息中心，它才真正落到界面上；界面反过来也证伪了最初的理由。这不是事后补的理由，而是这次比对暴露出的结论：\textbf{「有没有界面」这条判据本身是模糊的}，它取决于界面定义到哪一层，用它做排除条件时要格外小心。

\subsection{U-02：下单链路上模块交互的推导}

\textbf{使用场景。} 模块划分完成后，需要说清模块之间怎么交互，这是组件图与时序图的依据。计划用「走一遍具体业务流程」的办法推导依赖关系，而不是直接罗列「A 依赖 B」。

\textbf{生成结果。} 按步骤走完「顾客提交订单 $\rightarrow$ 订单变待支付、库存被占住」这条链路，得出主控为预订与订单，依赖账号与权限、定价、库存、审计，并说明门店与房源、支付与退款不参与的原因；据此提出三根单向箭头。

\textbf{优缺点分析。} 长处是\textbf{按业务流程逐步推导，比直接罗列更易发现遗漏}，并且能主动指出「哪些模块不该参与」——这一步用「不参与」的方式把边界划清楚了。不足很严重：\textbf{它给出的具体例子本身是错的}。它举「M4 库存释放后调 M5 会成环」作为需要避免的例子，但库存释放由订单取消触发，本来就是 M5 调 M4，箭头方向没变，不存在环。真正成环的位置是 M5 与 M6 之间（下单创建支付单是 M5$\rightarrow$M6，支付回调改订单状态是 M6$\rightarrow$M5）。\textbf{而且它不会自己发现这个错误}，需要人拿需求原文逐条去对。

\textbf{修改优化后的结果。} 下单链路确定为 M5 主控、三根单向依赖（M5$\rightarrow$M3、M5$\rightarrow$M4、M5$\rightarrow$M1）；M5 与 M6 之间的环用\textbf{事件}拆开——M6 支付成功后发布「支付已成功」事件，M5 订阅后自行推进订单状态，M6 不调用 M5。同时确定两处细节：支付单在用户点「去支付」时才创建，不在下单时创建，使支付网关故障不波及下单；库存与订单在同一事务内\textbf{先扣库存再建订单}，因为扣库存是链路上唯一会失败的关键动作，放在最前可以让失败的请求尽早退出。这两条后来成了第 3 章的决策四与决策一。

\subsection{U-03：模块职责与接口的五处分歧判断}

\textbf{使用场景。} 按九个模块起草职责与接口表时，有五处存在两种都说得通的写法，需要定一种并落到表里：① 横切依赖（M1、M9 被所有模块依赖）在组件图上怎么画；② 排房归门店与房源还是房态与工单；③ 22:00 未到店判定归预订模块还是房态模块；④ 线上可售上限由店长直接设定还是由公式算出；⑤ 取消后释放库存用同步调用还是事件。

\textbf{生成结果。} 针对五处各给出甲乙两案，并附上自己的倾向。

\textbf{优缺点分析。} 长处是\textbf{先把「存在两种写法」这件事摆出来}，这比只给一个答案有用得多——有几处（例如排房到底归谁）在提问之前完全没想到要选；附上倾向也便于快速判断。不足是\textbf{它的倾向未必对，而且错得很有代表性}：第 ① 处它倾向「不画箭头，只在图下注一句『所有模块依赖 M1 与 M9』」，理由是画了会糊掉，却没意识到那样这两个模块根本不会出现在组件图上，而作业明确要求写明模块划分，缺了这两块等于交了一张残缺的图。\textbf{它优化的是图的观感，不是作业的达标线}。

\textbf{修改优化后的结果。} 逐条定为：

\begin{enumerate}[leftmargin=3.2em]
  \item \textbf{改为与 AI 相反}：把 M1、M9 画成底部一条「公共支撑」带，各业务模块汇一根总箭头过去。同时在正文里明写这个折中本身违反「画全」原则及原因，不让它悄悄消失。
  \item 采纳：M2 只管房间档案，占房与释放房号是 M7 的事。
  \item 采纳：M5 的定时任务自己判定 22:00 未到店，M7 只回答「这单今天入住没有」——改的是订单状态，就该由订单模块主导。
  \item 采纳：可售上限以公式（房间总数 $-$ 维修 $-$ 线下保留）为天花板，店长可在此基础上\textbf{手动下调，但不得高于公式值}。这一条同时满足 REQ-006 的「可设置」与 REQ-033 的「按公式重算」，也是第 4 章那条公式的由来。
  \item 采纳：M5 \textbf{同步调用} M4 释放库存，不经过事件。释放失败会让库存长期被占、直接少卖房，这种事必须当场知道。
\end{enumerate}

\subsection{U-04：19 屏跨屏数据一致性核对}

\textbf{使用场景。} 19 屏界面按屏逐个生成时，每一屏单独看都自洽；但界面之间共享同一批数字（房量、出租率、价格、营收），一旦某一屏单独改动，就会出现「两套事实」。因此在 19 屏齐备后，做了一轮逐屏对账。

\textbf{生成结果与发现的问题。} 核对的方法是：先固定 1.5 节的六条数据基准，再把每一屏上出现的数字抽出来与基准比对。发现的问题可分为三类：

\begin{itemize}
  \item \textbf{跨屏不一致}：片区概览的营收环比写作 $+4.8\%$，但片区合计营收 ¥113,721 与前一日推算值对应的是 $+2.0\%$；门店排班页的「今日当班」人数与下方排班表实际当班行数对不上；顾客端各屏的门店评分写作 4.8，而后台门店档案与详情页统一是 4.6。
  \item \textbf{单屏金额算错}：订单列表里「一居室公寓 2 晚」的已付金额写作 ¥588，应为 ¥588 $\times$ 2 = ¥1,176；酒店详情页的「剩 3 间」与两晚库存取最小值 $\min(4,1)=1$ 矛盾，应为「剩 1 间」。
  \item \textbf{规则被界面悄悄改写}：支付页的取消政策原文暗示「18:00 之后可自助取消」，与 REQ-021（18:00 后转前台代提交 + 店长审批）直接冲突；平台端的退款案例时间落在 18:00 之后却走的是自助链路。这类错误最危险——\textbf{数字错了只是难看，规则错了会让人以为系统支持一个它并不支持的操作}。
  \item \textbf{补屏时另起一套事实}（补齐 17--19 屏时才暴露）：门店与人员屏的门店清单是另编的六家（富阳、临平、武林门等），房量与出租率与 15 屏完全不同，页头的「今日出租率 78.6\%」也不等于任一组按房量加权的值。原因是\textbf{后生成的屏没有去读已有屏的数字}，它只从「六家门店」这个描述出发重新编了一遍。同一批里还有身份冲突：同一个人在一部分屏是「湖滨银泰店长」、在平台三屏却是「片区运营」，审计屏甚至出现「操作人 $\times\times$ · 批准人 $\times\times$」的同人自批——正好撞在 REQ-030 的回避规则上。
\end{itemize}

\textbf{优缺点分析。} 长处是\textbf{生成速度与一致的单屏质量}：19 屏在结构、间距、状态色上高度一致，这是逐屏手画很难达到的。不足是\textbf{它没有跨屏的全局状态}——每次只在一个上下文中生成一屏，看不见其他屏上的数字；而设计里最需要一致的东西恰恰是跨屏的。此外它对业务规则的理解是「读到的这句话」，而不是「这条规则在所有界面上应该怎么体现」，所以会在不同屏上给出互相矛盾的倾向。

\textbf{修改优化后的结果。} 建立「单一事实源」：48 间房的六种房态相加必须等于 48、48 间 $\times$ 价格 $\times$ 出租率必须能解释当日营收、片区 6 店按房量加权的出租率必须等于 72.7\%。三道校验写进报告 1.5 节，之后任何一屏改数字都要过这三道；上列问题逐条修正并重新导出。

第三类问题（补屏另起一套事实）的修法与数字不同，它要求\textbf{先指定一屏为权威副本}：门店清单以 15 屏为准，19 屏逐行照抄同一组门店名、房量与出租率，房源结构的两条汇总也按 296 间重新配平，使两行各自的和都等于门店房量之和；身份冲突则改成一个人只有一个角色——片区运营另有其人，店长名只出现在本店语境里，审计屏的自批条目改写为「区间内无需审批」，删除申请的申请人与复核人拆成两个不同的人。这类修改必须\textbf{改到源文件再重新导出}，不能在导出图上修图，否则下一轮重导又回到错误状态。

\subsection{U-05：报告自身的图表与引用自检}

\textbf{使用场景。} 报告正文较长，图、表、决策条目、章节之间大量互相引用。写完后需要确认「文中提到的编号确实存在」。

\textbf{生成结果与发现的问题。} 逐章交叉检查引用，发现两处：

\begin{itemize}
  \item 第 2 章追溯矩阵把 REQ-015 引向「3.5 决策五」、REQ-016 引向「决策六」，但当时的决策表只有六行，且第五条讲的是状态机、第六条讲的是价格快照——\textbf{编号对得上，内容对不上}。
  \item 第 1 章与第 5 章开头都写「画了 14 张 UML 图」，实际目录里是 16 张；类图说明写「23 个类」，实际数是 24 个。
\end{itemize}

\textbf{优缺点分析。} 长处是\textbf{一致性检查适合交给模型}：这类问题是纯机械的比对，不涉及业务判断，AI 做得又快又全，比人工翻页可靠。不足是\textbf{它会在改一处时引入新的一处}——决策表扩成八条之后，如果没有重新对一遍所有引用，就会出现「新加的第七条没人引用、旧的第五条被错指」的二次问题。\textbf{一致性不是一次检查，而是每次改动后都要重跑的东西}。

\textbf{修改优化后的结果。} 决策表由六条扩为八条（新增决策七「查询走缓存、写落主库」与决策八「鉴权在网关一次做完」），并把第 2 章的引用改为决策七 / 决策八 / 决策三；图的张数与类的数量按实际改。这里有一个值得记下的教训：\textbf{写「第 N 章」「决策 N」这类相对编号，比写「3.5 节决策三」更容易在扩写后失效}，因此定稿时应统一按「章-节-序号」的绝对位置引用。

\subsection{大模型如何辅助软件设计}

把五次记录横过来看，AI 在本设计中的实际作用是清楚的：

\begin{table}[H]
\centering
\caption{AI 在设计各环节的实际作用与人的介入点}
\tabfont
\begin{tabularx}{\textwidth}{@{}L{2.9cm}L{4.8cm}X@{}}
\toprule
设计环节 & AI 做了什么事 & 人必须做的事 \\
\midrule
需求理解 & 把 36 条需求批量归类到模块 & 判断分类是否违背交付约束；发现规模失衡 \\
模块划分 & 给出多种划分维度与各自代价 & 按作业要求排除「无界面」的模块；定最终切法 \\
依赖推导 & 按业务流程逐步推依赖 & 拿需求原文核对每一个具体例子 \\
图与界面生成 & 生成 PlantUML 与 19 屏界面，风格统一 & 检查规则是否被界面悄悄改写；跨屏数字对账 \\
文本与编号 & 起草正文、维护图表编号与引用 & 判断编号背后指的到底是不是那个意思 \\
\bottomrule
\end{tabularx}
\end{table}

一句话概括：\textbf{AI 擅长的是「把已有约束铺开」，不擅长的是「判断哪条约束在这件事上更硬」。} 它能快速生成十个候选、九条依赖、十九屏界面，但它不知道「每个模块要有界面」是这个作业的达标线、「18:00 后不能自助取消」是业务红线——它只知道这些字出现在输入里。

\subsubsection{能否设计一套工作流，让人只做检查校对}

作业的结尾问题值得认真回答。结论是：\textbf{能做到「人只做检查校对」，但校对这一步不能简化成「看一眼」。}

一条可行的工作流分五步，前四步可以自动化，第五步不行：

\begin{enumerate}[leftmargin=3.2em]
  \item \textbf{结构化输入}。把需求写成带编号的条目（本设计的 REQ-001～036），把数据基准写成一张表（本设计的 1.5 节），把交付约束单独列成清单（图的类别与下限、每模块需有界面、提交物结构）。这一步是整条流水线的价值来源：\textbf{能被机器校验的东西，前提是它先被写成了可校验的形状}。
  \item \textbf{分单元生成}。按图、按屏、按章分别生成，每个单元附一句「本单元落了哪几条 REQ」。
  \item \textbf{机器校验}。三类自动检查：PlantUML 语法编译通过；REQ 覆盖矩阵与需求清单做差集，找出没被任何单元引用的需求；跨单元数值比对，把同一实体在不同单元上的数字列出来看是否相等。
  \item \textbf{差异回写}。把第三步的结果原样喂回生成环节，让模型只改被指出的位置，而不是整段重写——避免 U-05 里那类「改一处引入另一处」。
  \item \textbf{人工校对}。人只做三件机器做不了的事：\textbf{判断约束的优先级}（作业要求与图面观感冲突时，哪一个是达标线——U-03 的第 ① 处）；\textbf{判断「算得对」之外是否「说得对」}（U-04 里「18:00 后自助取消」数字全对但规则错了，这类机器查不出来）；\textbf{判断需求本身是否成立}（REQ-035、REQ-036 至今标为待确认，任何生成物都不该替实验一把它确认掉）。
\end{enumerate}

因此答案是：\textbf{可以让一个人一天之内完成原本一周的设计工作量，但那个人的角色从「画图的人」变成了「定判据的人」}。他的产出不再是一张图，而是那三条机器能读的输入——需求编号、数据基准、交付约束清单。这三样东西写不清，整条流水线产出的就是一堆漂亮而互相矛盾的东西；写清了，检查校对才真的只是检查校对。

反过来说，本设计里 AI 出的每一个错，追根到底都不是「模型能力不够」，而是\textbf{判据没有提前写下来}：模块划分时「有没有界面」是个模糊判据，于是它把通知塞进了订单模块；数据基准是逐屏才浮现的，于是出现了 ¥588 与 ¥1,176；决策编号是相对的，于是扩写后引用全错。\textbf{把判据前置，是这个工作流唯一真正的难点，也是最不能被 AI 代劳的部分。}
